\documentclass[11pt]{article}

\setlength{\parindent}{0pt}
\usepackage{xltxtra}
\usepackage{hyperref}
\hypersetup{hidelinks}
\usepackage{url}
\urlstyle{tt}
\usepackage{xcolor}
\definecolor{CVBlue}{RGB}{23,110,191}
\usepackage{calc}
\usepackage{graphicx}
\usepackage{tikz}
\usetikzlibrary{calc}
\usepackage{fontspec}
\usepackage{xeCJK}
\usepackage{enumitem}
\CJKsetecglue{} %% 取消中文与数字之间的间隙

%% 主文档字体设置
\setmainfont[
    Path = fonts/Main/,
    Extension = .otf,
    BoldFont = texgyretermes-bold.otf, % 加粗字体
]{texgyretermes-regular.otf} % 正文字体

% 中文字体设置
\setCJKmainfont[
    Path = fonts/hansans/,
    Extension = .ttf,
    BoldFont = NotoSansSC-Bold.ttf, % 加粗字体
]{NotoSansSC-Regular.otf} % 正文字体

\usepackage{relsize}
\usepackage{xspace}

% 使用 fontawesome（部分图标）
\usepackage{fontawesome} 

% A4纸，上下左右边距
\usepackage[
    a4paper,
    left=1.2cm,
    right=1.2cm,
    top=1.5cm,
    bottom=1cm,
    nohead
]{geometry}

\renewcommand{\baselinestretch}{1.5} % 行间距设为1.5

\usepackage{titlesec}
\usepackage{enumitem}
\setlist{noitemsep} % 取消列表项间的额外间距
%\setlist{nosep} % 取消所有垂直间距
\setlist[itemize]{topsep=0.25em, leftmargin=*}
\setlist[enumerate]{topsep=0.25em, leftmargin=*}

% --- 用于控制【不同项目之间】的垂直距离 ---
\newlength{\interProjectSpacing}
\setlength{\interProjectSpacing}{0.9em} % <--- 在此调整项目之间的距离
\newcommand{\projectsep}{\vspace{\interProjectSpacing}}

% --- 用于控制【项目标题】与下方【项目描述】的距离 ---
\newlength{\intraProjectTitleSep}
\setlength{\intraProjectTitleSep}{0.4em} % <--- 在此调整标题和描述的距离
\newcommand{\titlebreak}{\\[\intraProjectTitleSep]}

% --- 用于控制【项目描述】与下方【要点列表】的距离 ---
\newlength{\intraProjectListTopSep}
\setlength{\intraProjectListTopSep}{0.2em} % <--- 在此调整描述和列表的距离

% =======================================================================

\titleformat{\section}         % 定制 \section 命令 
{\large\bfseries\raggedright} % 将 section 标题设置为大号、粗体且左对齐
{}{0em}                      % 可用于为所有 section 添加前缀（如“章节...”）
{}                           % 可用于在标题前插入代码
[{\color{CVBlue}\titlerule}]  % 在标题后插入一条横线
\titlespacing*{\section}{0cm}{*1.6}{*1.2}


\begin{document}
\pagenumbering{gobble}

%%%% 利用tikz来定位照片 (如需使用个人照片，请确保路径下有 avatar.jpg)
\begin{tikzpicture}[remember picture, overlay] 
    \node[anchor = north east] at ($(current page.north east)+(-2cm,-1.2cm)$) {\includegraphics[height=3cm]{avatar.jpg}};
  \end{tikzpicture}%
  %%%% 利用tikz来定位学校Logo (如需使用浙大Logo，请确保路径下有 zju.png)
  \begin{tikzpicture}[remember picture, overlay] 
    \node[anchor = north west] at ($(current page.north west)+(0.5cm,+1.0cm)$) {\includegraphics[height=6cm]{zju.png}};
  \end{tikzpicture}%
\centerline{\LARGE\bfseries{邵楷城}} 

\centerline{\normalsize{\faPhone\ 13567814497 \quad \faEnvelopeO\ \href{mailto:1350843130@qq.com}{1350843130@qq.com} \quad \faMapMarker\ 浙江省宁波市}} 
    
\section{\makebox[\widthof{\faGraduationCap}][c]{\color{CVBlue}\faGraduationCap}\ 教育背景}    
\textbf{浙江大学} \hfill 2024.09 -- 至今\\[0.5em] 
电子信息（人工智能方向） \quad 硕士（在读） \quad 交互数据分析团队 ZJUIDG（隶属计算机辅助设计与图形系统 (CAD\&CG) 全国重点实验室）
\begin{itemize}[nosep]
    \item \textbf{获奖情况：浙江大学优秀研究生，浙江大学五好研究生，一等助学金}
\end{itemize}

\projectsep

\textbf{杭州电子科技大学} \hfill 2020.09 -- 2024.06\\[0.5em] 
计算机科学与技术 \quad 本科 
\section{\makebox[\widthof{\faUsers}][c]{\color{CVBlue}\faUsers}\ 项目经历}

% --- 第一个项目 ---
\textbf{大唐-光伏电站多源融合系统 (校企合作)} \hfill 2024.07 -- 2024.11 \titlebreak
项目描述：与大唐集团合作，围绕光伏时序大数据展开，涵盖数据治理、异常识别、故障诊断和可视化平台开发四个核心方向，旨在提升光伏电站的管理与诊断效率。
\begin{itemize}[nosep, topsep=\intraProjectListTopSep]
    \item \textbf{模型训练与优化：} 负责光伏电站图纸数据的预处理与模型构建，应用多种深度学习算法（结合时序与图像特征的多模态分析）识别组串异常，显著提升模型故障诊断性能。
    \item \textbf{数据治理协助：} 配合团队采用分层架构和自适应算法进行数据清洗，提升原始数据质量。
    \item \textbf{方案撰写：} 独立负责并完成该合作项目投标书的技术方案撰写与梳理工作。
\end{itemize}

\projectsep

% --- 第二个项目 ---
\textbf{基于多层注意力机制的多图融合交通流量预测方法} \hfill 2024.01 -- 2024.03 \titlebreak
项目描述：本科优秀毕业设计项目。针对复杂交通网络场景，提出一种基于多层注意力机制与多图融合的创新预测框架，实现高精度的交通流量预测。
\begin{itemize}[nosep, topsep=\intraProjectListTopSep]
    \item \textbf{算法创新与实现：} 提出一种计算道路功能相似性的新算法，构建并融合两种不同关系的图网络；在聚合信息后，同步提取短时序与全局时序信息。
    \item \textbf{模型复现与调优：} 深入调研并学习图卷积神经网络 (GCN) 相关前沿文献，熟练运用 PyTorch 框架独立完成深度学习模型的搭建、训练、复现与效果完善。
    \item \textbf{全流程实践：} 完整掌握并实践了从数据获取、图网络构建、多层注意力机制设计到模型评估的深度学习开发全流程。
\end{itemize}

\section{\makebox[\widthof{\faCogs}][c]{\color{CVBlue}\faCogs}\ 技能特长}
\begin{itemize}[nosep]
    \item \textbf{编程语言：} 熟练掌握 Python、C/C++
    \item \textbf{算法与框架：} 熟悉 PyTorch 等主流深度学习框架；对图卷积神经网络 (GCN)、注意力机制及多模态特征融合有深入理解与实战经验。
    \item \textbf{英语能力：} CET-6 (504分)；CET-4 (529分)；考研英语一 (88分)，具备流畅的英文学术文献阅读与技术文档查阅能力。
\end{itemize}

\end{document}
